\begin{figure*}[t]
\centering%
\begin{subfigure}[t]{0.35\textwidth}
\centering%
%\includegraphics[width=\textwidth]{schema/model}
%\includegraphics[width=\textwidth]{schema/model-recon}
\includegraphics[width=\textwidth]{schema/model-recon-gray}
\caption{从经验中学习动力学}
\label{fig:dynamics}
\end{subfigure}%
\hfil%
\begin{subfigure}[t]{0.35\textwidth}
\centering%
\includegraphics[width=\textwidth]{schema/imagination}
\caption{在想象中学习行为}
\label{fig:behavior}
\end{subfigure}%
\hfil%
\begin{subfigure}[t]{0.245\textwidth}
\centering%
\includegraphics[width=\textwidth]{schema/interaction}
\caption{在环境中行动}
\label{fig:interaction}
\end{subfigure}%
\caption[]{Dreamer 的组成部分。(a) 基于过去经验构成的数据集，智能体学习将观测和动作编码为紧凑的潜在状态（\iconstate{}），例如通过重建来学习，并预测环境奖励（\iconreward{}）。(b) 在紧凑潜在空间中，Dreamer 预测状态价值（\iconvalue{}）和动作（\iconaction{}），并通过想象轨迹反向传播梯度，使未来价值预测最大化。(c) 智能体编码当前 episode 的历史，计算当前模型状态，并预测下一步将在环境中执行的动作。智能体伪代码见 \cref{alg:agent}。}
\label{fig:schema}
\end{figure*}
